\begin{frame}{Background 1: few-shot and chain-of-thought prompting}
\begin{block}{Definition}
Both techniques provide the LLM a pre-prompt \(\mathbf{p}\) which appended to the question \(\mathbf{x}\) is supposed to achieve a more precise output \(\mathbf{y}\). \\
In particular \(\mathbf{p}\) provides a representation of \(\mathbf{k}\) question-answer-tuples \(\left\{(\mathbf{x}_i, \mathbf{y}_i)\right\}_{k\in \mathds{N}}\)
which are supposed to tell the LLM how to derive his solutions.
\end{block}
\begin{table}
\small
\begin{tabular}{@{}p{0.45\textwidth} | p{0.5\textwidth}@{}}
\toprule
\textbf{few-shot} & \textbf{CoT} \\
\midrule
\(\mathbf{p} \equiv \langle \mathbf{x}_1 \cdot \mathbf{y}_1 \rangle \concat \langle \mathbf{x}_2 \cdot \mathbf{y}_2 \rangle \concat \dots \concat \langle \mathbf{x}_k \cdot \mathbf{y}_k \rangle\) & \(\mathbf{p} \equiv \langle \mathbf{x}_1 \cdot \mathbf{t}_1 \cdot \mathbf{y}_1 \rangle \concat \langle \mathbf{x}_2 \cdot \mathbf{t}_2 \cdot \mathbf{y}_2 \rangle \concat \dots \concat \langle \mathbf{x}_k \cdot \mathbf{t}_k \cdot \mathbf{y}_k \rangle\)\\
Then  \[\mathbf{p} \concat \mathbf{x} \underset{\text{LLM}}{\rightarrow} \mathbf{y}\] & Then \[\mathbf{p} \concat \mathbf{x} \underset{\text{LLM}}{\rightarrow} \mathbf{t} \cdot \mathbf{y}\] 
%\bottomrule
\end{tabular}
\end{table}
Where \(\cdot\) and \(\concat\) are appropriate separators in order to represent \(\left\{(\mathbf{x}_i, \mathbf{y}_i)\right\}_{k\in \mathds{N}}\).\\
\textrightarrow\ \CoT\ extends \textit{few-shot} by providing an \textbf{additional derivation} \(\mathbf{t}_i\). \\
Then \(\mathbf{p}\) is a representation of \(\left\{(\mathbf{x}_i, \mathbf{t}_i, \mathbf{y}_i)\right\}_{k\in \mathds{N}}\).

\end{frame}